%% Formule de l'operateur : Stationnarite_Covar_XY
%% Dans ce fichier, les �ventuels accents sont cod�s
%% en ISO-8859-15 (aka latin9)
%% La ligne suivante donne trois exemples 
%% e accent aigu : �, e accent grave : �, e-dans-l'o : �

$M$ est le nombre de sous-p�riodes sur lesquelles on calcule la covariance. Donc $M$ vaut $T_2 /T_1$ ( $T_2$ > $T_1$ et $T_2$ multiple de $T_1$ )\\

On d�finit : \\
$N_1$ : nombre d'�chantillons dans la p�riode $T_1$ \\

La covariance sur la $i$�me p�riode $T_1$ est calcul�e par :
$$
(\overline{X'Y'})_i = \frac{1}{N_1}\left[ \sum_j X_j.Y_j - \frac{1}{N_1} \sum_j X_j \sum_j Y_j \right] 
$$
Puis la covariance moyenne vaut
$$
\overline{X'Y'}_{moy} = \frac{1}{M} \sum_i (\overline{X'Y'})_i
$$
L'indice de stationnarit� vaut :
$$
Istat = \left\vert  \frac{ \overline{ X'Y'}  - \overline{X'Y'}_{moy} }
                   { \overline{ X'Y'} } \right\vert
$$
avec $\overline{ X'Y'}$ : covariance calcul�e sur toute la p�riode $T_2$